\begin{abstract}
When an adversary observes the last hop to a third-party tool, or is the tool, an LLM Agent's privacy cover cannot be discarded before it reaches the service: decoys must execute as real calls, consuming quota, touching state and possibly re-entering the planner's context. We characterize when cover must execute, classify executed cover by termination, state, feedback and dispatcher, and propose an evaluation that binds residual inference, task completion, business-state integrity and cost to the same protected executions. In 3,000 single-host executions with two fixed-weight planners over MCP and TLS~1.3/HTTP/2, controller cover that skips the model wait is removable: given call boundaries, label-free timing removal raises recognition from 75 to 112 (Qwen) and from 80 to 108 (Llama) of 120 tasks. History-matched dispatch removes this shortcut, yet adapted attacks still recognize 93--106/120, a significant gain in three of four cells; without call boundaries, the prespecified removal attack loses its advantage. Against a shadow planner, the strongest tested attack reaches but does not exceed a pair-ambiguity floor we derive (7/12 for one shadow), which only more shadows or label-aware sampling can lower. In a stress control, goal completion misses collateral state damage in 95 of 102 completed tasks. Existing Agent defenses report an adapted attacker and matched task utility in no common setting; applying our protocol with declared tolerances shows that the policy to deploy depends on the declared observer.
\end{abstract}
